\documentclass[12pt,letterpaper]{article}

\usepackage{amsmath,amssymb,amsthm}
\usepackage{mathtools}
\usepackage{geometry}
\usepackage{hyperref}
\usepackage{booktabs}
\usepackage{array}
\usepackage{xcolor}
\usepackage{microtype}
\usepackage{enumitem}
\usepackage{mdframed}
\usepackage{tcolorbox}
\usepackage[numbers,sort&compress]{natbib}

\geometry{margin=1in}

% ── Theorem environments ──────────────────────────────────────────────────────
\theoremstyle{plain}
\newtheorem{theorem}{Theorem}[section]
\newtheorem{lemma}[theorem]{Lemma}
\theoremstyle{definition}
\newtheorem{definition}[theorem]{Definition}
\theoremstyle{remark}
\newtheorem{remark}[theorem]{Remark}

% ── Colored boxes ─────────────────────────────────────────────────────────────
\tcbuselibrary{breakable,skins}

\newtcolorbox{flagbox}{
  breakable, enhanced,
  colback=red!4, colframe=red!60!black,
  fonttitle=\bfseries\small, title=Expert Flag,
  boxrule=0.8pt, left=4pt, right=4pt, top=4pt, bottom=4pt
}

\newtcolorbox{improvebox}{
  breakable, enhanced,
  colback=blue!4, colframe=blue!60!black,
  fonttitle=\bfseries\small, title=Proposed Improvement,
  boxrule=0.8pt, left=4pt, right=4pt, top=4pt, bottom=4pt
}

\newtcolorbox{opbox}{
  breakable, enhanced,
  colback=green!4, colframe=green!60!black,
  fonttitle=\bfseries\small, title=Open Problem,
  boxrule=0.8pt, left=4pt, right=4pt, top=4pt, bottom=4pt
}

% ── Custom commands ───────────────────────────────────────────────────────────
\newcommand{\kB}{k_{\mathrm{B}}}
\newcommand{\lP}{\ell_{P}}
\newcommand{\kt}{\tilde{\kappa}}
\newcommand{\kteff}{\tilde{\kappa}_{\mathrm{eff}}}
\newcommand{\Sent}{S_{\mathrm{ent}}}
\newcommand{\Ment}{M_{\mathrm{ent}}}
\newcommand{\asc}{\alpha_{\mathrm{screen}}}
\newcommand{\tr}{\operatorname{Tr}}
\newcommand{\sint}{S_{\mathrm{int}}}

% ─────────────────────────────────────────────────────────────────────────────
\title{\textbf{Expert Analysis: Journal White Paper on Entanglement Entropy Density as a Gravitational Source}\\[6pt]
\large Tasks 1--5: Framework Stress Test, Theoretical Tightening, OP3 Project Plan,\\
Experimental Sharpening, and Publication Strategy}

\author{Expert Physics Review\\
\small Responding to: K.\ Monette, \textit{Entanglement Entropy Density as a Gravitational Source},\\
\small Draft March 19, 2026}

\date{March 2026}

\begin{document}
\maketitle

\tableofcontents
\newpage

% ─────────────────────────────────────────────────────────────────────────────
\section*{Task 1: Framework Restatement and Stress Test}
\addcontentsline{toc}{section}{Task 1: Framework Restatement and Stress Test}
\label{sec:task1}

The white paper introduces a phenomenological extension of general relativity in which the internal entanglement structure of a laboratory quantum device contributes to the gravitational stress-energy tensor. The coupling is mediated by a coarse-grained scalar field $\phi(x) \equiv \Sent(x)$ representing the density of entanglement between subsystems of the device (bits per cubic meter), treated as an emergent, effective degree of freedom rather than a new fundamental field. This quantity is operationally defined: it is the von Neumann entropy $\sint(t;A) = -\kB\tr_A[\rho_A\ln\rho_A]$ of a chosen bipartition $A|B$ of the device, divided by $\kB\ln 2$ and by the cell volume. Crucially, the paper correctly identifies that only candidate~(3) --- the internal inter-subsystem entanglement entropy --- has the right behavior: it \emph{decreases} under decoherence (as $\rho_D$ approaches a product state over subsystems), unlike total von Neumann entropy or system--environment entropy, which both increase.

The coupling to gravity is introduced through a minimal effective action containing a kinetic term for $\phi$ and a linear potential $\alpha\phi$, ensuring that variation with respect to $g^{\mu\nu}$ produces a stress-energy tensor of the cosmological-constant type: $T^{\mathrm{ent}}_{\mu\nu} = \Lambda_{\mathrm{ent}}(x)g_{\mu\nu}$ with $\Lambda_{\mathrm{ent}} = \tilde{\kappa}c^4\Sent/(8\pi G\kB\ln 2)$. The paper uses the action to argue that total stress-energy is covariantly conserved: the scalar field equation of motion $\beta\Box\phi = \alpha$ ensures $\nabla^\mu T^{\mathrm{total}}_{\mu\nu} = 0$ on-shell. The Bianchi identity then implies that ordinary matter is \emph{not} independently conserved where $\nabla\Sent \neq 0$ --- the framework predicts energy-momentum exchange between matter and the entanglement sector. The paper correctly notes that the specific estimate $\kt = -1/4$ is a motivated ansatz under a Planck-scale holographic regulator hypothesis and is not derivable from first principles; $\kt$ is treated as a free phenomenological parameter throughout.

For the radial structure, the paper correctly shows that the entanglement-induced acceleration correction scales as $\Delta a(R) \propto \Ment(<\!R)/R^2$ for any bounded source, resolving the earlier confusion about $1/R$ vs.\ $1/R^2$. The experimental program is appropriately modest: NISQ-scale devices with 50 qubits are the near-term target; macroscopic BEC interferometry is demoted due to extreme spatial screening. The falsification criterion is reformulated as a bound on the observable combination $|\kteff|\,\Delta\Ment/R_{\mathrm{eff}}^2$ rather than a direct bound on $\kt$, which is the correct framing given that $\asc$ (the ratio of actual to ideal time-integrated entanglement) is currently unknown.

\textbf{What the framework does not claim:} It does not derive $\kt$ from first principles. It does not assert that the scalar field $\phi$ is fundamental. It does not claim that existing equivalence-principle experiments (MICROSCOPE, Asenbaum et al.) directly constrain $\kt$, since their test masses have negligible and nearly identical $\Sent$. It does not provide the Lindblad derivation of $\asc$ (this is explicitly flagged as an open problem). The paper is admirable in these omissions.

% ─────────────────────────────────────────────────────────────────────────────
\section{Task 2: Theoretical Tightening}
\label{sec:task2}

\subsection{Consistency check of the covariant formulation}

The stress-energy definition and conservation statements are consistent with standard GR conventions. The decomposition $T^{\mathrm{total}}_{\mu\nu} = T^{\mathrm{matter}}_{\mu\nu} + T^{\mathrm{ent}}_{\mu\nu}$ is well-posed. The metric signature is not stated in the white paper; this should be fixed as $(-,+,+,+)$ in the final manuscript to remove ambiguity for readers coming from different conventions. Specifying $G_{\mu\nu} = R_{\mu\nu} - \tfrac{1}{2}Rg_{\mu\nu}$ explicitly in the definition of the modified field equation is also advisable.

Three implicit assumptions about $\Sent(x)$ need to be stated explicitly:

\begin{enumerate}
\item \textbf{Smoothness.} The coarse-graining procedure in Section~2.3 produces a piecewise-constant field on cells of volume $\Delta V$. The paper's Poisson equation analysis assumes $\rho_{\mathrm{ent}}(x)$ is smooth. A sentence should state: ``We treat $\Sent(x)$ as smooth at scales large compared to the device lattice spacing; this is the effective-field-theory approximation appropriate below the Planck scale.''

\item \textbf{Bipartition dependence.} The definition $\sint(t;A)$ depends on the choice of bipartition $A|B$. The paper does not specify how this choice is made across spacetime or whether a canonical choice exists. For a uniform linear chain of qubits, the natural choice is the half-chain bipartition; for a 2D array, a geometric half-plane. This must be stated explicitly, otherwise $\Sent(x)$ is not well-defined as a scalar field.

\item \textbf{Locality.} The coarse-grained field $\Sent(x)$ is assigned to a spatial cell, but entanglement entropy is a nonlocal quantity (it depends on correlations across the cut, not just within the cell). The paper should note that the coarse-graining breaks translational invariance in a hardware-dependent way and that $\Sent(x)$ is therefore not a local observable in the usual QFT sense --- it is a \emph{collective, coarse-grained} quantity.
\end{enumerate}

\subsection{The critical tension in the effective action}

\begin{flagbox}
The paper's effective action has a fundamental tension that must be addressed before publication.

The scalar field $\phi(x) \equiv \Sent(x)$ is simultaneously:
\begin{enumerate}[label=(\alph*)]
\item \textbf{Externally prescribed:} $\phi(x)$ is determined by the quantum state of the device, prepared and controlled by experimenters. It is \emph{not} a freely propagating degree of freedom.
\item \textbf{Dynamical:} The action's variation with respect to $\phi$ yields the field equation $\beta\Box\phi = \alpha$, which is an autonomous dynamics for $\phi$ independent of the quantum state.
\end{enumerate}
These are incompatible. If $\phi$ is externally prescribed (as it must be for experimental control), then varying the action with respect to $\phi$ is not a valid operation --- the field equation is not satisfied by the experimentally prepared state.

Specifically: for a cluster state of 50 qubits, $\Sent(x)$ is a specific function of space determined by the quantum circuit. This function does not satisfy $\beta\Box\Sent = \alpha$ in general. The action's $\phi$-variation is therefore fictitious.

\textbf{Two resolutions exist:}

\textbf{Resolution A (Constrained action).} Treat $\phi$ as a \emph{fixed external source} (not dynamical), and add it to the action only through the coupling $\alpha\phi\sqrt{-g}$. Then the field equation is dropped, and $T^{\mathrm{ent}}_{\mu\nu}$ is computed from the variation with respect to $g^{\mu\nu}$ only, giving simply $T^{\mathrm{ent}}_{\mu\nu} = \alpha\phi g_{\mu\nu} - \tfrac{\beta}{2}(\nabla\phi)^2 g_{\mu\nu} + \beta\nabla_\mu\phi\nabla_\nu\phi$. Conservation is ensured not by the $\phi$ equation of motion but by the identity $\nabla^\mu G_{\mu\nu} = 0$ together with the specific $\phi$ profile determined by the quantum state.

\textbf{Resolution B (Source term).} Promote $\phi$ to a fully dynamical field sourced by the quantum device: add a source term $J(x)\phi$ to the action, where $J(x)$ encodes how the quantum device drives the entanglement sector. Then $\phi$ satisfies $\beta\Box\phi = \alpha + J(x)$, and the quantum state enters through $J(x)$. In this case, $\phi$ deviates from $\Sent^{\mathrm{quantum}}$ by the propagator of the Klein-Gordon equation --- an interesting physical prediction, but also a significant complication.

\textbf{Recommended:} Resolution~A. It preserves the original physical picture (experimentally controlled $\phi$), avoids the fictitious field equation, and keeps the framework minimal. The kinetic term should be retained because it contributes to the stress-energy and gives $T^{\mathrm{ent}}_{\mu\nu}$ a physically richer structure than pure $\Lambda_{\mathrm{ent}}g_{\mu\nu}$.
\end{flagbox}

\subsection{Proposed improvement to the effective action}

\begin{improvebox}
\textbf{Replacement text for Section 3.3 of the white paper.}

To provide an action-based derivation of $T^{\mathrm{ent}}_{\mu\nu}$ without treating $\phi$ as a free dynamical field, we adopt the following approach. We introduce $\phi(x) = \Sent(x)$ as an \emph{external scalar source field} determined by the quantum state of the device --- it is not varied in the action. The effective action for the gravitational sector, in the presence of the fixed source $\phi(x)$, is:
\begin{equation}
  S[g;\phi]
  = \int d^4x\sqrt{-g}
    \left[
      \frac{c^3}{16\pi G}R
      - \frac{\beta}{2}(\nabla\phi)^2
      - \alpha\phi
    \right],
  \label{eq:action-fixed}
\end{equation}
where $\alpha \equiv \tilde{\kappa}c^4/(8\pi G\kB\ln 2)$ and $\beta > 0$ is a coupling with dimensions of (energy density) $\times$ (length/entropy$)^2$. Varying with respect to $g^{\mu\nu}$ alone gives the Einstein equations:
\begin{equation}
  G_{\mu\nu}
  = \frac{8\pi G}{c^4}\left[
    T^{\mathrm{matter}}_{\mu\nu}
    + \beta\!\left(\nabla_\mu\phi\nabla_\nu\phi
    - \tfrac{1}{2}g_{\mu\nu}(\nabla\phi)^2\right)
    + \alpha\phi\,g_{\mu\nu}
  \right].
  \label{eq:EFE-improved}
\end{equation}
The entanglement stress-energy is:
\begin{equation}
  T^{\mathrm{ent}}_{\mu\nu}
  = \beta\!\left(\nabla_\mu\phi\nabla_\nu\phi
  - \tfrac{1}{2}g_{\mu\nu}(\nabla\phi)^2\right)
  + \alpha\phi\,g_{\mu\nu}.
  \label{eq:Tent-improved}
\end{equation}
Conservation of total stress-energy, $\nabla^\mu T^{\mathrm{total}}_{\mu\nu} = 0$, is guaranteed by $\nabla^\mu G_{\mu\nu} \equiv 0$ and does not require any field equation for $\phi$. Instead, it imposes a consistency condition:
\begin{equation}
  \nabla^\mu T^{\mathrm{matter}}_{\mu\nu}
  = -\nabla^\mu T^{\mathrm{ent}}_{\mu\nu}
  = -\beta\,(\Box\phi)\,\nabla_\nu\phi
    + \beta\,\nabla_\nu\!\left(\tfrac{1}{2}(\nabla\phi)^2\right)
    - \alpha\,\nabla_\nu\phi.
  \label{eq:matter-conservation}
\end{equation}
The right-hand side is a specific function of the externally prescribed $\phi(x)$ profile. It predicts an anomalous force on matter wherever $\nabla\Sent \neq 0$, with magnitude proportional to $\alpha\,|\nabla\Sent|$. This is a concrete, composition-independent fifth-force prediction that can be bounded experimentally.

In the limit $\nabla\phi \approx 0$ (slowly varying entanglement density), Eq.~\eqref{eq:Tent-improved} reduces to $T^{\mathrm{ent}}_{\mu\nu} \approx \alpha\phi\,g_{\mu\nu} = \Lambda_{\mathrm{ent}}g_{\mu\nu}$, recovering the cosmological-constant ansatz as the leading approximation. The kinetic term $\beta(\nabla\phi)^2$ provides the next correction, physically encoding the energy cost of entanglement-density gradients.
\end{improvebox}

\subsection{Why a potential $V(\phi)$ and curvature coupling are not needed}

The reviewer prompt asks whether a potential $V(\phi)$ or curvature coupling $\xi R\phi^2$ should be added. The answer is: not in the first paper. Both would introduce additional free parameters without adding falsifiable content at NISQ scales. A curvature coupling $\xi R\phi$ (linear in $\phi$) would be degenerate with the $\alpha\phi$ term in Eq.~\eqref{eq:action-fixed} at leading order in perturbation theory. A nonlinear potential $V(\phi) = m^2\phi^2/2$ would give $\phi$ a mass $m$, causing the entanglement field to be Yukawa-suppressed at distances $R \gg \hbar/(mc)$. Unless there is a physical reason to expect such a mass, it should not be introduced. These are appropriate additions for the technical backup paper (Paper~II) as model extensions to be constrained by experiment.

% ─────────────────────────────────────────────────────────────────────────────
\section{Task 3: OP3 --- Full Lindblad Project Plan}
\label{sec:task3}

\subsection{Platform choice and justification}

We choose a \textbf{fixed-frequency superconducting transmon chain} (IBM Quantum architecture) for the following reasons:
\begin{itemize}
\item Hamiltonian and Lindblad parameters are publicly available from hardware calibration data
\item Gate set tomography (GST) data exists for characterizing Kraus maps
\item Tensor network methods (MPO-TDVP) are well-developed for 1D chains up to $N \sim 100$
\item Circuit depth can be varied continuously and systematically
\end{itemize}

\subsection{Precise definitions for the OP3 program}

\textbf{System Hamiltonian} (in rotating frame, two-level approximation):
\begin{equation}
  H_{\mathrm{sys}}
  = \sum_{k=1}^{N}\frac{\omega_k}{2}\sigma^z_k
  + \sum_{k=1}^{N-1}J_k\!\left(\sigma^+_k\sigma^-_{k+1} + \mathrm{h.c.}\right),
  \label{eq:H-transmon}
\end{equation}
with $\omega_k/2\pi \in [4.5, 5.0]$\,GHz (sampled from published IBM calibration) and $J_k/2\pi = 5$\,MHz (nearest-neighbor exchange coupling via bus resonator).

\textbf{Lindblad operators and rates:}
\begin{align}
  L_{1,k} &= \sqrt{\gamma_1}\,\sigma^-_k,
  & \gamma_1 &= 1/T_1 \approx 2\pi\times5\,\text{kHz} \quad (T_1 = 200\,\mu\text{s}),
  \label{eq:L1} \\
  L_{\phi,k} &= \sqrt{\gamma_\phi/2}\,\sigma^z_k,
  & \gamma_\phi &= 1/T_\phi \approx 2\pi\times8\,\text{kHz} \quad (T_\phi = 125\,\mu\text{s}).
  \label{eq:Lphi}
\end{align}
These are state-of-the-art parameters from IBM Falcon-R10 and Eagle processors (2024--2025).

\textbf{Master equation:}
\begin{equation}
  \dot{\rho}
  = -\frac{i}{\hbar}[H_{\mathrm{sys}},\rho]
  + \sum_{k=1}^{N}\sum_{\alpha\in\{1,\phi\}}
    \!\left(L_{\alpha,k}\rho L_{\alpha,k}^\dagger
    - \tfrac{1}{2}\{L_{\alpha,k}^\dagger L_{\alpha,k},\rho\}\right).
  \label{eq:lindblad}
\end{equation}

\textbf{Definition of $M_{\mathrm{ent}}(t)$ for this platform:}

For a chain of $N$ qubits, we use the half-chain bipartition $A = \{1,\ldots,N/2\}$, $B = \{N/2+1,\ldots,N\}$:
\begin{equation}
  M_{\mathrm{ent}}(t) = \sint(t; A) = -\kB\,\tr_A[\rho_A(t)\ln\rho_A(t)],
  \quad \rho_A(t) = \tr_B\rho(t).
  \label{eq:Ment-def}
\end{equation}

\textbf{Screening factor:}
\begin{equation}
  \asc(N,\tau) = \frac{\int_0^\tau M_{\mathrm{ent}}(t)\,dt}{\int_0^\tau M_{\mathrm{ent}}^{\mathrm{ideal}}(t)\,dt},
  \label{eq:asc-def}
\end{equation}
where $M_{\mathrm{ent}}^{\mathrm{ideal}}(t)$ is computed with $\gamma_1 = \gamma_\phi = 0$ (ideal unitary evolution).

\textbf{Initial states for the program:}

\begin{itemize}
\item \textbf{Cluster state (1D):} $|\mathrm{Cl}\rangle = \prod_{\langle jk\rangle} \mathrm{CZ}_{jk}\,|+\rangle^{\otimes N}$.
  Ideal half-chain entropy: $M_{\mathrm{ent}}^{\mathrm{ideal}} = \kB\ln 2$ (maximal for any bipartition of a 1D cluster state).
\item \textbf{GHZ state:} $|\mathrm{GHZ}\rangle = (|0\rangle^{\otimes N} + |1\rangle^{\otimes N})/\sqrt{2}$.
  Ideal half-chain entropy: $M_{\mathrm{ent}}^{\mathrm{ideal}} = \kB\ln 2$ (same magnitude, but different decay behavior).
\item \textbf{Product state (Branch B):} $|+\rangle^{\otimes N}$.
  Ideal entropy: $M_{\mathrm{ent}}^{\mathrm{ideal}} = 0$.
\end{itemize}

\subsection{Staged numerical research plan}

\begin{opbox}
\textbf{Stage 1 (Months 1--2): $N \leq 10$, exact diagonalization.}

\textbf{Method:} Full Lindblad master equation~\eqref{eq:lindblad} evolved via matrix exponential (Expokit or SciPy \texttt{scipy.linalg.expm}) on the $4^N$-dimensional Liouvillian. For $N = 10$: Liouvillian dimension $= 4^{10} \approx 10^6$; feasible on a single workstation.

\textbf{Computations:}
\begin{enumerate}
\item For initial states $|\mathrm{Cl}\rangle$ and $|\mathrm{GHZ}\rangle$ on $N \in \{2, 4, 6, 8, 10\}$ qubits:
  \begin{enumerate}[label=(\alph*)]
  \item Evolve $\rho(t)$ under~\eqref{eq:lindblad} for $0 \leq t \leq 5\max(T_1, T_\phi)$.
  \item Compute $M_{\mathrm{ent}}(t)$ via~\eqref{eq:Ment-def} at each time step.
  \item Compute $\asc(N,\tau)$ via~\eqref{eq:asc-def} for $\tau = T_\phi, 2T_\phi, 5T_\phi$.
  \end{enumerate}
\item Also compute the \emph{site-resolved} entropy $s_k(t) = -\kB\tr_k[\rho_k(t)\ln\rho_k(t)]$
  to understand spatial structure of entanglement decay.
\end{enumerate}

\textbf{Output:} Plots of $M_{\mathrm{ent}}(t)$ and $\asc(N,\tau)$ vs.\ $N$ and $\tau$ for both cluster and GHZ states. Comparison of decay rates. Identification of which state gives larger $\asc$ at fixed $N, \tau$.

\textbf{Expected finding:} Cluster states should show $\asc \sim e^{-\gamma_\phi\tau}$ (single-qubit decay time $T_\phi$) while GHZ states show $\asc \sim e^{-N\gamma_\phi\tau}$ (collective decay time $T_\phi/N$). Quantitative verification of this prediction is the primary output of Stage~1.
\end{opbox}

\begin{opbox}
\textbf{Stage 2 (Months 2--4): $N \leq 20$, gate-error Kraus maps.}

\textbf{Method:} Replace ideal unitary gates in the cluster state preparation circuit with their Kraus-map representations, obtained from gate set tomography (GST) data. GST data is publicly available for IBM devices from published calibration reports.

\textbf{Additional modeling:} For each two-qubit CZ gate on qubits $(j,j+1)$, replace $U_{\mathrm{CZ}}$ with a Kraus map $\mathcal{E}_{\mathrm{CZ}}^{(j)}$ characterized by a $4\times 4$ process matrix $\chi^{(j)}$. The gate error rate $\varepsilon_{2Q} \approx 5\times10^{-3}$ (IBM Falcon, 2025) enters as:
\begin{equation}
  \mathcal{E}_{\mathrm{CZ}}(\rho) = (1-\varepsilon_{2Q})U_{\mathrm{CZ}}\rho U_{\mathrm{CZ}}^\dagger
  + \varepsilon_{2Q}\,\mathcal{D}_{\mathrm{depol}}(\rho),
  \label{eq:kraus}
\end{equation}
where $\mathcal{D}_{\mathrm{depol}}$ is a local depolarizing map.

\textbf{Computations:}
\begin{enumerate}
\item For $N \in \{10, 15, 20\}$ qubits, prepare the cluster state circuit using Kraus maps.
\item Evolve under~\eqref{eq:lindblad} with the gate errors applied at each gate layer.
\item Compute $M_{\mathrm{ent}}(t)$, $\asc(N,\tau)$, and the \emph{total entanglement budget}:
\begin{equation}
  \mathcal{E}_{\mathrm{total}}(N,d) \equiv \int_0^\tau \sum_k s_k(t)\,dt,
\end{equation}
where $d$ is the circuit depth. Optimize $d$ to maximize $\mathcal{E}_{\mathrm{total}}$.
\end{enumerate}

\textbf{Output:} $\asc(N,\tau)$ as a function of $N$, $\tau$, and circuit depth $d$. A fit of the form $\asc \approx A\,e^{-N^\alpha\gamma_\phi\tau - d\varepsilon_{2Q}}$ to identify scaling exponents $\alpha$ and the relative contributions of decoherence vs.\ gate errors. The value of $\alpha$ determines whether $\asc$ scales polynomially (cluster state, $\alpha \approx 0$) or exponentially (GHZ, $\alpha = 1$) in $N$.
\end{opbox}

\begin{opbox}
\textbf{Stage 3 (Months 4--8): $N = 50$, tensor network methods.}

\textbf{Method:} For $N = 50$, the Liouvillian Hilbert space dimension is $4^{50} \approx 10^{30}$, far beyond exact diagonalization. Use matrix product operator (MPO) methods:

\begin{enumerate}
\item Represent $\rho(t)$ as an MPO with bond dimension $\chi$.
\item Evolve under the Lindblad equation using the time-dependent variational principle (TDVP) for MPOs (implemented in TeNPy, ITensor, or quimb).
\item Monitor entanglement entropy of $\rho(t)$ and $\rho_A(t)$ at each time step.
\end{enumerate}

\textbf{Bond dimension requirements:} For cluster states under 1D Lindblad dynamics, the operator entanglement entropy of $\rho(t)$ grows logarithmically in time, suggesting that MPO bond dimension $\chi \sim 10^2$--$10^3$ should suffice for $N = 50$ up to $t \sim T_\phi$. For GHZ states, the rapid entanglement collapse makes the dynamics easier (lower bond dimension after decoherence begins).

\textbf{Computations:}
\begin{enumerate}
\item Prepare the $N = 50$ cluster state as an MPS with bond dimension $\chi_0 = 2$ (exact for the cluster state).
\item Convert to MPO representation for the Lindblad evolution.
\item Evolve for $0 \leq t \leq 3T_\phi$; compute $M_{\mathrm{ent}}(t)$ and $\asc(50,\tau)$.
\item Repeat for product state (Branch B) as control.
\end{enumerate}

\textbf{Output:} Full $\asc(N,\tau)$ curve for $N = 50$. Extrapolation to the $N\to\infty$ limit if the scaling law from Stage~2 holds. A definitive answer to the question: is $\asc$ nonzero and significant for a 50-qubit cluster state on a transmon platform with realistic noise?

\textbf{The pivotal result from Stage 3 is:} If $\asc(50, T_\phi) \gtrsim 10^{-2}$, the experimental program is viable. If $\asc \lesssim 10^{-4}$, the NISQ experiment will not produce a useful bound and the framework's near-term experimental program must be reconsidered.
\end{opbox}

\begin{opbox}
\textbf{Stage 4 (Months 8--12): Analytic bounds and scaling theory.}

Using the results of Stages 1--3, derive or conjecture:
\begin{enumerate}
\item A \emph{lower bound} on $\asc$ for cluster states using the quantum data processing inequality:
  $\sint(t; A) \geq \sint(0; A) - \gamma_\phi\,\xi\,t$ for some correlation length $\xi$,
  where the inequality follows from the fact that local dephasing can only reduce entanglement at a rate bounded by the dephasing rate times the correlation length.
\item A \emph{critical $N$} at which $\asc$ drops below $1\%$ of its ideal value.
  From the Stage 2 scaling law $\asc \approx A\,e^{-\gamma_\phi\tau}$ for cluster states:
  $N^* \sim |\ln(0.01)|/\gamma_\phi\tau = $ a computable number.
\item The \emph{optimal circuit}: the cluster state preparation circuit $C^*$ that maximizes $\mathcal{E}_{\mathrm{total}}(N,d)$ for fixed $N$ and hardware parameters. This is an optimization problem over the space of native-gate sequences.
\end{enumerate}

\textbf{Output:} A table of $(N, d^*, \asc^*)$ for the optimal cluster state circuits on the IBM transmon platform. A scaling law $\asc(N) \approx f(N)$ with analytic form. Plots showing the tradeoff between $N$ (larger $\Ment$) and $\asc$ (smaller due to decoherence), with the optimal operating point $N^*$.
\end{opbox}

\subsection{How OP3 outputs plug back into the white paper}

Once $\asc(N,\tau)$ is computed, the following replacements can be made in the white paper:

\begin{itemize}
\item Section~5.1: Replace ``$\asc$ is an open parameter'' with ``For a 50-qubit transmon cluster state with $T_1 = 200\,\mu$s, $T_2 = 100\,\mu$s, and $\tau = T_\phi/2$, we find $\asc = [numerical value] \pm [error]$ from MPO-Lindblad simulation.''
\item Section~6: Replace the bound on $|\kteff|\,\Delta\Ment/R_{\mathrm{eff}}^2$ with an explicit numerical value using the computed $\asc$ and measured $\Delta\Ment$.
\item Paper~III (experimental proposal): Include the $\asc$ calculation as a theoretical prediction that the experiment should test. A significant deviation of measured $\asc$ from the Lindblad prediction would itself be a result --- it would indicate either new physics or an error in the open-system model.
\end{itemize}

% ─────────────────────────────────────────────────────────────────────────────
\section{Task 4: Experimental Proposal Sharpening}
\label{sec:task4}

\subsection{Fully specified experimental configuration}

\textbf{Platform:} $6\times 6 = 36$ qubit $^{87}$Rb neutral-atom array in reconfigurable optical tweezers. Lattice spacing: $a = 5\,\mu$m. Two-qubit gates via Rydberg blockade ($|70S_{1/2}\rangle$, blockade radius $r_b = 8\,\mu$m for nearest-neighbor coupling only).

\textbf{Qubit basis:}
$|0\rangle = |5S_{1/2}, F=1, m_F=0\rangle$,
$|1\rangle = |5S_{1/2}, F=2, m_F=0\rangle$
(clock transition, $\omega_{\mathrm{hf}}/2\pi = 6.835$\,GHz).

\textbf{Branch A (high entanglement): 2D cluster state.}
\begin{enumerate}
\item Initialize: $|+\rangle^{\otimes 36}$ via $\pi/2$ pulses on clock transition.
\item Apply CZ gates to all nearest-neighbor horizontal pairs (one layer, $\tau_{\mathrm{CZ}} = 200$\,ns per gate, parallelized: total time $200$\,ns).
\item Apply CZ gates to all nearest-neighbor vertical pairs (one layer, $200$\,ns).
\item Result: 2D cluster state. Half-array entanglement entropy $M_{\mathrm{ent}}^A \approx 6\kB\ln 2$ (for a $3\times 6$ vs $3\times 6$ bipartition).
\item Total preparation time: $\tau_{\mathrm{prep}} \approx 1\,\mu$s.
\end{enumerate}

\textbf{Branch B (low entanglement): Matched product state.}
\begin{enumerate}
\item Initialize: $|+\rangle^{\otimes 36}$ (same as Branch A).
\item Apply CZ gates to all horizontal pairs with \emph{zero interaction time} (Rydberg laser on but detuned by $\Delta_{\mathrm{off}} = 10\Omega$, giving gate fidelity $\approx (1-\Omega^2/\Delta_{\mathrm{off}}^2) = 99\%$ for the identity operation).
\item Apply the same ``null'' vertical CZ gates.
\item Result: Product state $|+\rangle^{\otimes 36}$ with $M_{\mathrm{ent}}^B = 0$.
\end{enumerate}

\textbf{Why this Branch A/B pair:} The 2D cluster state and the $|+\rangle^{\otimes 36}$ product state have identical single-site density matrices ($\rho_k = \mathbf{I}/2$ for every qubit $k$) and identical single-site energies $\langle H_k\rangle = 0$ for the clock Hamiltonian. The two-body correlators $\langle\sigma^z_j\sigma^z_k\rangle$ differ between cluster state and product state, but for nearest-neighbor Ising-type Hamiltonians, the energy difference is:
\begin{equation}
  \Delta E_{AB} = \langle H_{\mathrm{sys}}\rangle_A - \langle H_{\mathrm{sys}}\rangle_B
  = \sum_{\langle jk\rangle} J_{jk}\underbrace{\left(\langle\sigma^+_j\sigma^-_k\rangle_A - \langle\sigma^+_j\sigma^-_k\rangle_B\right)}_{\text{same for both (both}=0\text{)}}.
  \label{eq:energy-diff}
\end{equation}
For the clock qubit Hamiltonian with no $XX$ or $YY$ coupling (only $ZZ$ from the Rydberg interaction during gates, zero after gates), both states have the same mean energy after state preparation. This is the crucial property ensuring that the \emph{mean-field} frequency shift is matched between branches.

\subsection{Timing and observable}

\textbf{Ramsey measurement protocol:}
\begin{enumerate}
\item Prepare Branch A (or B) state via the protocol above.
\item Apply global $\pi/2$ pulse on clock transition: $R_x(\pi/2)$.
\item Free precession for $T_{\mathrm{Ramsey}} \in \{1, 10, 100\}\,\text{ms}$.
\item Apply global $\pi/2$ pulse: $R_x(\pi/2)$ or $R_x(-\pi/2)$.
\item Measure population in $|1\rangle$ via site-resolved fluorescence.
\item Repeat $n_{\mathrm{shots}} = 10^4$ times.
\end{enumerate}

\textbf{Differential phase:}
\begin{equation}
  \Delta\phi_{AB}
  = \phi^A - \phi^B
  = (\omega^A_{\mathrm{eff}} - \omega^B_{\mathrm{eff}})\,T_{\mathrm{Ramsey}},
  \label{eq:Delta-phi}
\end{equation}
where $\omega^{A,B}_{\mathrm{eff}}$ is the effective transition frequency in each branch.

\textbf{Signal formula:} The entanglement coupling predicts a frequency shift:
\begin{equation}
  \frac{\delta\omega_{\mathrm{hf}}}{\omega_{\mathrm{hf}}}
  = \frac{\delta g_{00}}{2}
  = -\frac{\Phi_{\mathrm{ent}}}{c^2},
  \label{eq:freq-shift}
\end{equation}
where $\Phi_{\mathrm{ent}}$ is the effective gravitational redshift potential from the entanglement source. In the Newtonian limit, from Section~4:
\begin{equation}
  \Phi_{\mathrm{ent}} = -\frac{\kappa_{\mathrm{eff}}\,\Ment}{R_{\mathrm{eff}}},
  \label{eq:Phi-ent}
\end{equation}
where $R_{\mathrm{eff}} \sim a = 5\,\mu$m is the characteristic distance scale within the array and $\Ment = M_{\mathrm{ent}}^A - M_{\mathrm{ent}}^B = 6\kB\ln 2$ (the differential entanglement). The measurable quantity translates to a bound:
\begin{equation}
  \boxed{
  |\kteff|\,\frac{\Delta\Ment}{R_{\mathrm{eff}}^2}
  < \frac{|\Delta\phi_{AB}|_{\mathrm{sens}}\,8\pi G\kB\ln 2}{c^2\,\omega_{\mathrm{hf}}\,T_{\mathrm{Ramsey}}}
  \cdot\frac{1}{a}.
  }
  \label{eq:bound}
\end{equation}
Substituting: $|\Delta\phi|_{\mathrm{sens}} = 10^{-3}$\,rad (quantum projection noise at $N=36$ qubits with $10^4$ shots), $T_{\mathrm{Ramsey}} = 0.1$\,s, $a = 5\times10^{-6}$\,m, $\omega_{\mathrm{hf}} = 2\pi\times6.835\times10^9$\,s$^{-1}$:
\begin{equation}
  |\kteff|\,\frac{\Delta\Ment}{a^2}
  < \frac{10^{-3}\times8\pi\times6.67\times10^{-11}\times9.57\times10^{-24}}
         {(3\times10^8)^2\times2\pi\times6.835\times10^9\times0.1}
  \times\frac{1}{5\times10^{-6}}
  \approx 3\times10^{-52}\;\text{kg}^{-1}\,\text{m}\,\text{s}^2.
  \label{eq:bound-numerical}
\end{equation}

\subsection{Systematic error budget: fatal vs.\ non-fatal}

\begin{table}[h]
\centering\small
\caption{Systematic error budget for the 36-qubit neutral-atom Ramsey experiment.}
\label{tab:systematics}
\begin{tabular}{@{}p{4cm}p{3cm}p{4cm}p{2cm}@{}}
\toprule
Systematic & Origin & Mitigation & Fatal? \\
\midrule
Differential Zeeman shift & Magnetic field $B(t)$ fluctuations & Clock transition ($m_F=0$) insensitive to first order; active shielding to $\delta B < 1$\,nG & \textbf{No} \\
\addlinespace
Differential AC Stark (clock laser) & Laser intensity $I(t)$ fluctuations & Common-mode rejection: same laser for A and B; interleave shots & \textbf{No} \\
\addlinespace
Differential two-body correlator frequency shift & $\langle\sigma^z_j\sigma^z_k\rangle_A \neq \langle\sigma^z_j\sigma^z_k\rangle_B$ during Ramsey & Choose A,B states with matched $\langle\sigma^z_j\sigma^z_k\rangle$ (cluster vs.\ product both give zero off-resonant Ising coupling in free precession) & \textbf{No} (for chosen state pair) \\
\addlinespace
Rydberg admixture during free precession & Rydberg laser not fully off & Verify zero admixture via spectroscopy; separate Rydberg pulse from Ramsey & \textbf{No} \\
\addlinespace
Differential SPAM errors & State-dependent loss during readout & Symmetrize A/B preparation schedule; verify SPAM via reference states & \textbf{No} \\
\addlinespace
Differential mean-field shift from density-density interactions & $n_{\mathrm{Rb}}^A \neq n_{\mathrm{Rb}}^B$ & Same trap configuration, same atom number; post-select on identical atom count & \textbf{Yes} (if different $N_{\mathrm{atoms}}$) \\
\addlinespace
Motional state difference (entangled vs.\ product) & Different mean kinetic energy for A vs.\ B & Verify temperature equivalence via sideband spectroscopy; Doppler-free clock transition & \textbf{No} \\
\addlinespace
Entanglement verification failure & SPAM + readout errors prevent confirming entanglement depth & Use entanglement witnesses with SPAM correction; require $W < 0$ with $>5\sigma$ confidence before including run & \textbf{Yes} (if unverified) \\
\bottomrule
\end{tabular}
\end{table}

\textbf{The two fatal systematics in detail:}

\textit{Fatal Systematic 1 --- Atom number mismatch.} If Branch A and Branch B have different mean atom numbers (e.g., due to different preparation-induced loss), then the density-density interaction shifts $(\delta\omega) \propto a^{(3)}_{\mathrm{Rb}}\langle n^2\rangle$ differ between branches, producing a spurious differential phase. Mitigation: post-select on runs where both branches have the same atom count (verified by site-resolved imaging). This reduces statistics by a factor of $\sim$ (fraction of runs with equal count), but does not reduce sensitivity per run.

\textit{Fatal Systematic 2 --- Unverified entanglement.} If the claimed cluster state is actually a product state (due to gate failures), the differential entanglement $\Delta\Ment = 0$ and any measured phase shift would be misattributed to the coupling. Mitigation: run entanglement witness measurements before every Ramsey sequence. The witness operator $W_{\mathrm{Cl}} = \sum_k (\mathbf{I} - \sigma^z_k\sigma^z_{k+1})/2$ should satisfy $\langle W_{\mathrm{Cl}}\rangle < 0$ for any genuine cluster state. Runs failing this criterion are excluded.

\textbf{What a null result means.} If $\Delta\phi_{AB} = 0$ to within $|\Delta\phi|_{\mathrm{sens}} = 10^{-3}$\,rad after $10^4$ shots, this establishes the bound~\eqref{eq:bound-numerical}. This is a meaningful constraint on the observable combination $|\kteff|\,\Delta\Ment/R_{\mathrm{eff}}^2$ regardless of the value of $\asc$ or the Planck-scale renormalization of $\kt$. The paper correctly frames this as its primary experimental output.

% ─────────────────────────────────────────────────────────────────────────────
\section{Task 5: Publication Strategy and Gap Analysis}
\label{sec:task5}

\subsection{Paper I: Theory Letter (4--6 pages)}

\textbf{Target journal:} Physical Review Letters or Physical Review D (Letters format).

\textbf{Core claim:} A dimensionally consistent, covariantly conserved phenomenological extension of GR in which operationally defined internal entanglement entropy density contributes to the stress-energy tensor. The key structural results are: (a) definition of $\Sent$ as internal inter-subsystem entanglement entropy density (not total or system-environment entropy); (b) the minimal effective action~\eqref{eq:action-fixed} with $\phi$ as an external source field; (c) $\Delta a \propto \Ment/R^2$ for bounded sources; (d) the fifth-force prediction from Eq.~\eqref{eq:matter-conservation}; (e) the observable combination $|\kteff|\,\Delta\Ment/R_{\mathrm{eff}}^2$ as the experimental falsification target.

\textbf{What to keep from the white paper:} Sections 2, 3 (with improvements from Task 2), 4, and 6.1--6.2.

\textbf{What to cut:} The Lindblad model details (move to Paper~II). The macroscopic BEC discussion (keep only a brief note). The detailed experimental protocol (Paper~III).

\textbf{Minimal additional results before submission:}
\begin{itemize}
\item Resolve the action tension (Task 2, Resolution~A): one paragraph.
\item Add the fifth-force observable: one boxed equation.
\item State the metric signature explicitly.
\item Verify the Bose et al.\ (2023) citation.
\end{itemize}

\subsection{Paper II: Technical Backup (20--30 pages)}

\textbf{Target journal:} Physical Review D (full article) or Classical and Quantum Gravity.

\textbf{Core content:}
\begin{itemize}
\item Full derivation of the effective action and covariant conservation (Section 3 expanded)
\item Complete OP3 results (Stages 1--3 from Task 3): $\asc(N,\tau)$ for 50-qubit transmon cluster states
\item OP8: Bianchi identity consequences and fifth-force bounds
\item OP10: Planck-scale renormalization group power-counting (estimate whether $\kt$ at the Planck scale renormalizes to an observable value at laboratory scales)
\item Regime hierarchy table: $\asc$ for few-qubit (Regime I), mesoscopic (Regime II), macroscopic (Regime III)
\end{itemize}

\textbf{Minimal additional results before submission:}
\begin{itemize}
\item Stages 1--2 of the OP3 plan (Months 1--4): numerical Lindblad for $N \leq 20$.
\item A power-counting estimate of OP10: write $\kt_{\mathrm{lab}} = \kt_{\mathrm{Planck}} \times (E_{\mathrm{lab}}/E_P)^\alpha$ and estimate $\alpha$ from dimensional analysis.
\item An explicit calculation of the fifth-force bound: for a test mass near a 50-qubit chip, what is the anomalous acceleration predicted by Eq.~\eqref{eq:matter-conservation}?
\end{itemize}

\subsection{Paper III: Experimental Proposal (12--15 pages)}

\textbf{Target journal:} Physical Review A (experimental) or npj Quantum Information.

\textbf{Core content:}
\begin{itemize}
\item Full specification of the 36-qubit neutral-atom Ramsey experiment (Task 4)
\item Explicit formula~\eqref{eq:bound} relating measured $\Delta\phi_{AB}$ to the bound on $|\kteff|\,\Delta\Ment/R_{\mathrm{eff}}^2$
\item Complete systematic error budget (Table~\ref{tab:systematics})
\item The OP3 prediction for $\asc$ (from Stage 3 of Task 3) as a theoretical input to the experiment
\item A comparison section: what would it mean if $\asc^{\mathrm{measured}} \neq \asc^{\mathrm{Lindblad}}$?
\end{itemize}

\textbf{Minimal additional results before submission:}
\begin{itemize}
\item Stage 3 of the OP3 plan (MPO-Lindblad for $N = 50$): numerical prediction of $\asc(50,\tau)$.
\item Explicit calculation of the mean-field energy matching for the cluster vs.\ product state pair (verify $\Delta E_{AB} = 0$ for the chosen Hamiltonian).
\item An estimate of the quantum projection noise floor: $|\Delta\phi|_{\mathrm{sens}} \sim 1/\sqrt{N_{\mathrm{atoms}}\,n_{\mathrm{shots}}}$.
\end{itemize}

\subsection{Timeline}

\begin{center}\small
\begin{tabular}{@{}cll@{}}
\toprule
Months & Deliverable & Paper \\
\midrule
0--1  & Action tension resolved; $\Sent$ defined; fifth-force equation added & Paper~I \\
1--2  & Paper~I submitted & \\
1--4  & OP3 Stage 1--2: Lindblad numerics $N \leq 20$ & Paper~II \\
4--8  & OP3 Stage 3: MPO-Lindblad $N = 50$; OP10 power-counting & Paper~II \\
4--6  & Mean-field energy matching; noise floor estimate & Paper~III \\
8     & Paper~II and Paper~III submitted & \\
\bottomrule
\end{tabular}
\end{center}

% ─────────────────────────────────────────────────────────────────────────────
\section*{Summary: The Three Most Important Actions}
\addcontentsline{toc}{section}{Summary}

\begin{enumerate}
\item \textbf{Fix the action tension (one afternoon).} Change the effective action to treat $\phi = \Sent$ as an external source field (not varied), and note that conservation is guaranteed by $\nabla^\mu G_{\mu\nu} \equiv 0$, not by a $\phi$ equation of motion. Use the replacement text in Task~2. This removes a logical inconsistency and makes Paper~I submittable.

\item \textbf{Specify the bipartition (one sentence).} The scalar field $\Sent(x)$ depends on the choice of bipartition $A|B$. State explicitly that the canonical choice for a 2D array is a geometric half-plane cut, and that the results are insensitive to the exact cut location for translationally invariant systems. This makes $\Sent$ well-defined as a spacetime scalar field.

\item \textbf{Run Stage 1 of OP3 (two months).} Solve the Lindblad master equation numerically for $N = 10$ transmon qubits in a cluster state and compute $\asc(10, T_\phi)$. This single number determines whether the experimental program is viable at NISQ scales and whether Papers~II and~III are worth pursuing. If $\asc(10, T_\phi) \lesssim 10^{-6}$, the framework's near-term experimental program requires fundamental revision. If $\asc \gtrsim 10^{-2}$, the path forward is clear.
\end{enumerate}

% ─────────────────────────────────────────────────────────────────────────────
\bibliographystyle{unsrtnat}
\begin{thebibliography}{99}

\bibitem{Jacobson1995}
T.~Jacobson, \textit{Phys.\ Rev.\ Lett.}\ \textbf{75}, 1260 (1995).

\bibitem{Verlinde2017}
E.~Verlinde, \textit{SciPost Phys.}\ \textbf{2}, 016 (2017).

\bibitem{Asenbaum2020}
P.~Asenbaum \textit{et al.}, \textit{Phys.\ Rev.\ Lett.}\ \textbf{125}, 191101 (2020).

\bibitem{MICROSCOPE2022}
P.~T.~Touboul \textit{et al.}, \textit{Phys.\ Rev.\ Lett.}\ \textbf{129}, 121102 (2022).

\bibitem{Blume2017}
R.~Blume-Kohout \textit{et al.}, \textit{Nat.\ Commun.}\ \textbf{8}, 14485 (2017).

\bibitem{Haah2017}
J.~Haah \textit{et al.}, \textit{IEEE Trans.\ Inf.\ Theory}\ \textbf{63}, 5628 (2017).

\bibitem{TeNPy}
J.~Hauschild and F.~Pollmann,
``Efficient numerical simulations with Tensor Networks: Tensor Network Python (TeNPy),''
\textit{SciPost Phys.\ Lect.\ Notes}\ \textbf{5} (2018).

\bibitem{Saffman2016}
M.~Saffman, \textit{J.\ Phys.\ B}\ \textbf{49}, 202001 (2016).

\bibitem{Vidal2004}
G.~Vidal, \textit{Phys.\ Rev.\ Lett.}\ \textbf{93}, 040502 (2004).

\bibitem{Haegeman2016}
J.~Haegeman \textit{et al.}, \textit{Phys.\ Rev.\ B}\ \textbf{94}, 165116 (2016).
[TDVP for MPOs]

\end{thebibliography}

\end{document}
